% ROOT V137: ONLY c=1; keep \hbar and G dimensionful.
\documentclass[11pt]{article}
\usepackage{amsmath,amssymb,booktabs}
\begin{document}
\section*{ROOT V137: index, gauge representations, anomaly tests}
\[c=1,\qquad [\hbar]=ML,\qquad [G_4]=L/M,\qquad [G_{10}]=L^7/M.\]
\[K_6=\mathbb{CP}^{2}\times S^2,\qquad
c_1(L_{q,m})=(2q+1)h+2mu.\]
\[\operatorname{ind}D_{q,m}=\frac{q(q+1)}{2}\,m.\]
\[\operatorname{ind}D=3\implies (q,m)\in\{(1,3),(-2,3),(2,1),(-3,1)\}.\]
\[H^0(\mathbb{CP}^2,\mathcal O(1))\simeq\mathbf{3}\ \text{or}\ \overline{\mathbf 3},\qquad
H^0(\mathbb{CP}^1,\mathcal O(2))\simeq\mathbf{3}_{SU(2)}.\]
\[Q_L:(2,2)\mapsto(\mathbf3,\mathbf2),\quad
u^c_L:(1,1)\mapsto(\mathbf1,\mathbf1),\quad
u^c_L\text{ with }Y=0\text{ is input}.\]
\[\mathcal A_{SU(3)^3}=2-1-1=0.\]
\[\mathcal A_{SU(3)^2Y}=2\cdot\frac12\frac16+\frac12\Big(-\frac23+\frac13\Big)=0.\]
\[\mathcal A_{SU(2)^2Y}=3\cdot\frac12\frac16+\frac12\Big(-\frac12\Big)=0.\]
\[\mathcal A_{Y^3}=6\Big(\frac16\Big)^3+3\Big(-\frac23\Big)^3+3\Big(\frac13\Big)^3
+2\Big(-\frac12\Big)^3+1=0.\]
\[\mathcal A_{\mathrm{grav}^2Y}=6\Big(\frac16\Big)+3\Big(-\frac23\Big)+3\Big(\frac13\Big)+2\Big(-\frac12\Big)+1=0.\]
\[3+1=4\equiv0\pmod{2}\qquad\text{(Witten SU(2) test).}\]
\[Y(u^c_L)=-\frac23,\qquad Y(d^c_L)=\frac13,
\qquad(q,m)_{u^c_L}=(q,m)_{d^c_L}=(-3,1).\]
\[\therefore\ \nexists\ f:\mathbb Z^2\longrightarrow\mathbb Q\text{ with }Y=f(q,m)\text{ for the chosen template}.\]
\[I_{12}\propto\big[\widehat A(TX)\,e^{c_1(L)/2}\mathrm{ch}(E)\big]_{12}
\quad\text{(10D anomaly consistency required; not evaluated here).}\]
\end{document}
